\documentclass{article}
\usepackage[margin=2cm]{geometry}
\usepackage{amsmath}
\usepackage{graphicx}
\usepackage{booktabs}
\usepackage[utf8]{inputenc}
\usepackage{ragged2e}
\setlength{\emergencystretch}{3em}
\begin{document}
\sloppy

the solutions of the F(4) gauged supergravity theory being studied here can be uplifted to
an AdS6  S4 background of massive type IIA, our solutions should be captured by the D4-D8
brane framework. To identify the details of the relevant brane configuration, we first recall.
from section that the group which is gauged in the supergravity theory is SU(2)R  G+, where
G+ is the additional gauge group arising from the presence of vector multiplets. Indeed, the
presence of n vector fields A allows for the existence of a gauge group G+ of dimension
dimG+ = n. The gauge group G+ in the bulk corresponds to a flavor symmetry group
ENg+1 of the boundary SCFT . The RG-flow triggered by the gauge coupling breaks this
symmetry group to SO(2N)  U(1) in the IR. Deformation by the relevant mass parameters
will generically break SO(2N) further. For the solution studied in this paper, an SO(2)
symmetry survives, which suggests that a minimal choice for the dual field theory would be
one with N = 1 (i.e. a single D8 brane). However, even in this minimal case the enhanced
gauge group E2  SU(2)  U(1) of the conformal fixed point is found to have dimension
dim E2 = 4, which suggests that the holographic dual to such a theory should contain at
least four bulk vector multiplets. Fortunately, it is possible to embed our n = 1 solution
Setting the fields of the three additional vector multiplets to vanish then reproduces exactly
the solutions explored in this paper. In fact, such an embedding is possible for any value
behavior of all single-mass deformations of En,+1 theories for any N. As such, we will carry
out the localization calculation in section for generic N. We will find that for every choice
of 1  Nf  7, one obtains a good match between the analytic field theory expression and.
our previous numerical results. Having addressed the identification of flavor symmetries, it.
is natural to interpret the holographic solutions of this paper as dual to RG flows emanating
from the same UV SCFTs that were found to be the duals of pure Roman's supergravity.
The flow is driven by three relevant operators of dimension  = 3, 4, 3, in addition to the
gauge coupling deformation which brings the non-Lagrangian UV SCFT to an IR Yang-
Mills-matter theory. In the IR, the three relevant deformations are interpreted respectively
as a mass term for the hypermultiplet scalars, a mass term for the hypermultiplet fermions,
explicit form of these deformations is shown in . To support this interpretation, we now
calculate the free energy of the mass-deformed USp(2N) gauge theory and compare it to the
holographic result displayed in Figure . For the unfamiliar reader, we will first reproduce
the results of , where the USp(2N) theory without mass deformation was studied. The.
techniques used for the mass-deformed theory will be the same, and the new calculation is.
presented in section .


\subsection*{0.1Undeformed USp(2N) gauge theory}


In , localization techniques were used to find the perturbative partition function of  = 1

\end{document}
